\documentclass[12pt]{article}
\usepackage[T1]{fontenc}
\usepackage[utf8]{inputenc}
\usepackage{lmodern}
\usepackage{textcomp}
\usepackage{enumitem}
\usepackage{geometry}
\geometry{margin=1in}
\begin{document}
\begin{center}
Answer Key - Political Science - Graduate Level Assessment 20 \
\small (Answers are ordered exactly as in the question paper.)
\end{center}

\section*{Multiple Choice}
\begin{enumerate}[label=\arabic*.]
\item \textbf{Answer:} C
\\ \textit{Options:} A) bread and butter.; B) war and peace.; C) guns and butter.; D) bombs and books.
\\ \textit{Note:} The phrase "guns and butter" illustrates the economic trade-off between a nation's military spending (defense) and its domestic welfare expenditures (social spending), highlighting the opportunity cost associated with allocating resources to one over the other.
\item \textbf{Answer:} D
\\ \textit{Options:} A) A growing effort by governments to prevent the spread of weapons, especially of WMDs.; B) A growth in volume of the arms trade.; C) A change in the nature of the defence trade, linked to the changing nature of conflict.; D) All of these options.
\\ \textit{Note:} The correct answer is "All of these options" because the trends in the global defense trade since the end of the Cold War include increased privatization, the rise of emerging markets, and a shift towards joint military operations, all of which are interconnected developments reflecting the evolving geopolitical landscape.
\item \textbf{Answer:} B
\\ \textit{Options:} A) The threat posed by the Soviet Union; B) Domestic concerns of the US; C) Soviet ideology; D) None of the above
\\ \textit{Note:} Revisionists argue that the primary cause of the Cold War stemmed from the United States' domestic concerns, as American leaders prioritized internal stability and the containment of communism to address fears of economic decline, social unrest, and the spread of leftist ideologies within the country.
\item \textbf{Answer:} B
\\ \textit{Options:} A) Surgical strike.; B) Indiscriminate attack.; C) Smart bomb.; D) Precision target.
\\ \textit{Note:} The phrase "indiscriminate attack" suggests a lack of precision and accountability, which undermines the notion of legitimacy typically associated with the principles of the Revolution in Military Affairs (RMA) that emphasize targeted and ethical military operations.
\item \textbf{Answer:} A
\\ \textit{Options:} A) It damaged support for the US model of political economy and capitalism; B) It created anger at the United States for exaggerating the crisis; C) It increased support for American global leadership under President Obama; D) It reduced global use of the US dollar
\\ \textit{Note:} The 2008 financial crisis undermined confidence in the US model of political economy and capitalism, as it revealed significant systemic flaws and led to widespread economic instability, prompting both domestic and international observers to question the efficacy and reliability of American economic practices.
\end{enumerate}

\section*{True / False}
\begin{enumerate}[label=\arabic*.]
\setcounter{enumi}{5}
\item \textbf{Answer:} True
\\ \textit{Options:} A) True; B) False
\\ \textit{Note:} The statement is true because the General Agreement on Tariffs and Trade (GATT), the International Monetary Fund (IMF), and the World Bank are foundational institutions that collectively promote free trade, financial stability, and economic development within the framework of the liberal international economic order.
\item \textbf{Answer:} False
\\ \textit{Options:} A) True; B) False
\\ \textit{Note:} International policy encompasses a broad range of diplomatic, military, cultural, and environmental agreements and interactions between countries, not limited to economic agreements.
\item \textbf{Answer:} True
\\ \textit{Options:} A) True; B) False
\\ \textit{Note:} Motivational realism posits that states act based on a combination of self-interest, which includes the pursuit of wealth, the promotion of ideological beliefs, and the innate tendencies of human greed, thereby affirming that greedy states are indeed driven by these motivations.
\item \textbf{Answer:} False
\\ \textit{Options:} A) True; B) False
\\ \textit{Note:} Intergovernmental organizations often lack the necessary authority, resources, and political will to implement comprehensive monitoring mechanisms that can effectively identify and constrain abusive leaders, leading to their limited effectiveness in this regard.
\item \textbf{Answer:} True
\\ \textit{Options:} A) True; B) False
\\ \textit{Note:} A multipolar system is defined by the distribution of power among three or more significant states, which influences international relations and dynamics, making the statement true.
\end{enumerate}

\section*{Long Form Response}
\begin{enumerate}[label=\arabic*.]
\setcounter{enumi}{10}
\item \textbf{Answer:} In the Realist perspective of international relations, states are viewed as the primary actors, prioritizing their national interests and security in an anarchic global system where no central authority exists. This focus on state sovereignty leads to a competitive and often conflictual environment, as states seek to maximize their power relative to others. Realism posits that international politics is driven by the struggle for power, where military capability and strategic alliances play crucial roles. The implication of this perspective is that cooperation among states is often temporary and contingent upon shifting power dynamics, thereby emphasizing the importance of a pragmatic, often cynical view of global politics where moral considerations take a backseat to national interests. Consequently, Realism shapes the way states interact, leading to a focus on realpolitik and a reluctance to trust in international institutions or norms that might constrain their sovereignty.
\\ \textit{Note:} correct because it accurately captures the Realist view that states are the central actors in an anarchic international system, emphasizing the importance of national interests, security, and the competitive pursuit of power, which shapes global political dynamics.
\item \textbf{Answer:} Historical materialism (HM) faces significant criticisms, particularly regarding its reductionist approach to capitalism and class struggle. Critics argue that HM often portrays capitalism as fundamentally self-serving, neglecting the complex interplay of social, cultural, and political factors that also shape historical developments. Additionally, the tendency of some HM proponents to assert the existence of objective historical laws is debated, with many scholars questioning their validity and applicability across diverse contexts. This overemphasis on class struggle can lead to a narrow historical analysis, potentially overlooking other critical dimensions of human experience and agency. Ultimately, these criticisms challenge the robustness of HM as a comprehensive framework for understanding historical phenomena.
\\ \textit{Note:} correct because it identifies the reductionist nature of historical materialism's analysis of capitalism and class struggle, critiques its neglect of the multifaceted influences on historical developments, and questions the validity of its claims to objective historical laws, all of which are central to understanding the criticisms of HM.
\end{enumerate}

\vfill
\noindent\textit{End of Answer Key}
\end{document}